\section{Post-saturation successor: closure-carrying navigation}

Planning abstraction and refinement, counterexample-guided refinement, bidirectional representation migration, world-model replanning, and terminal-commitment models already provide strong navigation machinery. We therefore treat these as donor transforms. P7 asks when such a donor-valid transform also preserves scientific \emph{task closure}.

Let $T$ be a donor transform with native preservation predicate $D(T)$ and let
\[
o=(o_{tot},o_{amb},o_{front},o_{obj},o_{epoch})
\]
record, in the bounded successor instance, obligation totality, obligation unambiguity, material-frontier resolution, objective/question semantic continuity, and closure-epoch currency. Define
\[
\operatorname{Carries}(T,o)=D(T)\wedge o_{tot}\wedge o_{amb}\wedge o_{front}\wedge o_{obj}\wedge o_{epoch}.
\]
The five coordinates instantiate the theorem and are not claimed universally minimal.

\paragraph{Donor conservativity.}
Adding the closure carrier does not change the donor-native planning/refinement/round-trip/replanning/terminal-commitment verdict.

\paragraph{Closure separation.}
For every registered donor transform and each non-inert closure coordinate, two transformations can have identical donor-visible validity and different scientific-closure inheritance. Thus successful refinement, representation transport or terminal commitment remains useful without self-authorizing task-global closure.

\paragraph{Closure refinement.}
When a donor-valid transform fails closure transport because a set $S$ of registered closure coordinates is unresolved or incorrect, resolving all members of $S$ restores closure carrying under the exact contract while resolving a proper subset does not. This extends the counterexample-guided-refinement idea to the scientific completion contract itself.

\paragraph{Compositional closure transport.}
Two closure-carrying transforms compose scientifically only when the target obligation contract emitted by the first is exactly the source contract consumed by the second, or a registered bridge proves equivalence. Donor-visible composability alone is insufficient.

\paragraph{Ideal-product equivalence.}
A donor product supplied with the same closure coordinates, bridge rules and composition predicate agrees extensionally with P7. No centralization advantage is claimed.

\paragraph{Exhaustive bounded support.}
The registered model evaluates 320 donor-transform/closure states with zero donor-conservativity and ideal-product mismatches. It contains 25 minimal one-coordinate closure separations, 31 donor-product nonclosure countermodels, 155 exact full closure-refinement successes, 1,055 proper-subset failures, 25 heterogeneous composition successes and 25 bridge-mismatch composition countermodels. An independent implementation reproduces the canonical row digest \texttt{25f40385714adb15bca298a8cfd2b7fe2b28c96bfe462f6b60583be8f735b95f}. Those counts carry a multiplicity and are reported with it: neither \texttt{carries} nor \texttt{compose} takes a donor argument, so 320 is 64 counted five times, 25 separations is 5 counted five times, 155 and 1,055 are 31 and 211 counted five times, and the 25 composition successes and 25 bridge countermodels are one of each counted once per ordered donor pair. Only the 31 nonclosure countermodels are 31 distinct facts.

P7 therefore widens constructively: mature navigation mechanisms are absorbed and retained, while task-global scientific closure is carried as an explicit obligation object that can be selectively refined and compositionally transported. The result is bounded formal evidence, not deployed-agent superiority or universal open-world completeness.
